\documentclass[aspectratio=169,12pt]{beamer}
\usepackage[T1]{fontenc}
\usepackage[utf8]{inputenc}
\usepackage{lmodern}
\usepackage{babel}
\babelprovide[import,main]{french}
\usepackage{amsmath,amssymb,booktabs,graphicx}
\usetheme{Boadilla}
\usecolortheme{default}
\definecolor{mone}{RGB}{0,114,178}
\definecolor{mtwo}{RGB}{213,94,0}
\definecolor{mthree}{RGB}{0,130,100}
\setbeamercolor{structure}{fg=mone}
\setbeamertemplate{navigation symbols}{}
\setbeamertemplate{footline}{\hfill\insertframenumber/\inserttotalframenumber\hspace{0.4cm}\vspace{0.2cm}}
\setbeamertemplate{headline}{}
\title{Approximation neuronale et optimisation\\d'une décision thermique}
\subtitle{Valeur, sensibilité et qualité de décision\\sur un benchmark analytique}
\author{Mhamed Halhoul}
\institute{\small
Faculté Polydisciplinaire de Larache\\[2pt]
M2 en cours --- Master\\
Modélisation mathématique et applications et apprentissages\\[3pt]
Module : Algorithmes métaheuristiques\\
Enseignant du module : Pr. Jamal Daoudi}
\date{}
\newcommand{\Mone}{\texttt{M1\_16}}\newcommand{\Mtwo}{\texttt{M2\_16}}
\newcommand{\Mthree}{\texttt{M1\_32}}
\newcommand{\chart}[2]{\centering\includegraphics[width=\linewidth,height=#2\textheight,keepaspectratio]{../report/figures/#1.pdf}\par}
\begin{document}
\begin{frame}[plain]
\titlepage
\end{frame}

\begin{frame}{Apprentissage et décision physique}
\[
\theta\in\mathbb R^{49}:\ \min_\theta\mathcal L(\theta),
\qquad p\in[0.02,0.20]:\ \min_p J_\lambda(p).
\]
La convexité physique en $p$ ne rend pas la perte convexe en $\theta$.
\vspace{0.25cm}
\begin{itemize}
\item Association : l'erreur de sensibilité classe-t-elle mieux la mauvaise décision ?
\item Comparaison : la supervision valeurs et dérivées améliore-t-elle le score moyen, sous quels coûts ?
\end{itemize}
\end{frame}

\begin{frame}{Température et référence physique M0}
\[
T(p)=\frac{e^{-\pi^2p}}{\sqrt2}+0.2e^{-4\pi^2p},
\quad J_\lambda(p)=T(p)+\lambda p.
\]
Capteur $x=1/4$ de $u_t=u_{xx}$, deux modes, bords homogènes.
Toutes les variables sont sans dimension.
\[
J_\lambda'=T'+\lambda,\quad
0<m=T''(b)\leq J_\lambda''.
\]
M0 : $b$ si $\lambda\leq\lambda_b$, $a$ si $\lambda\geq\lambda_a$,
sinon dichotomie sur l'unique racine intérieure.
\[
\lambda_b=-T'(b)<\lambda_a=-T'(a).
\]
Référence flottante et bornes numériques conditionnelles.
\end{frame}

\begin{frame}{Réseau tanh à 49 paramètres}
\[
\theta=(w,b,v,d),\qquad 16+16+16+1=49.
\]
\[
z=\frac{2(p-a)}{b-a}-1,\qquad \alpha_N=\frac{100}{9}.
\]
\[
h_i=\tanh(w_iz+b_i),\qquad
\widehat T=d+\sum_i v_i h_i.
\]
\[
\widehat T'(p)=\alpha_N\sum_i v_iw_i(1-h_i^2).
\]
$c_{\rm out}=0$, $s_{\rm out}=1$. Dérivée par rapport à $p$ physique.
L'architecture n'impose ni monotonie ni convexité.
\end{frame}

\begin{frame}{Valeurs seules ou valeurs et dérivées}
\[
\mathcal L_1=\frac1n\sum_j\frac{(\widehat T-T)^2}{S_0^2},
\quad
\mathcal L_2=\mathcal L_1+\frac1n\sum_j
                    \frac{(\widehat T'-T')^2}{S_1^2}.
\]
\[
S_0=T(a)-T(b)\simeq0.57295,\qquad
S_1=-T'(a)\simeq9.31369.
\]
M1 apprend les valeurs. M2 ajoute les sensibilités.
\vspace{0.2cm}
Coefficient différentiel 1, aucun facteur $1/2$, aucune régularisation.
Échelles fixes issues de l'oracle, sans estimation sur diagnostic.
Gradients analytiques des 49 paramètres.
\end{frame}

\begin{frame}{Protocole apparié et coût des labels}
\centering
\begin{tabular}{lrrrr}
\toprule
Condition & Sites & $n_T$ & $n_{T'}$ & Coût si $q_1=1$\\
\midrule
\Mone &16&16&0&16\\
\Mtwo &16&16&16&32\\
\Mthree &32&32&0&32\\
\bottomrule
\end{tabular}
\par\vspace{0.3cm}
\raggedright
20 graines, mêmes paramètres initiaux par graine.
Adam full-batch : 3000 updates, snapshot 300 secondaire.
\par
$\eta=0.001$, $\beta_1=0.9$, $\beta_2=0.999$, $\epsilon=10^{-8}$ hors racine.
Compteur et moments conservés.
\par
257 sites diagnostiques disjoints. $q_1=1$ est conventionnel :
même coût de labels ne signifie pas mêmes sites ou même temps.
\end{frame}

\begin{frame}{Garanties classiques sous hypothèses}
\[
R\leq2\varepsilon+\alpha_{\rm opt}
\]
Erreur uniforme $\varepsilon$ et défaut global d'optimisation majorés.
Une RMSE ne fournit pas $\varepsilon$.
\[
|p-p^\star|\leq r/m,\qquad R\leq r^2/(2m).
\]
$r$ est le résidu KKT physique, avec les signes adaptés aux frontières.
\par\vspace{0.2cm}
Décision neuronale : meilleur des 1025 candidats, sans raffinement.
Aucune convexité du réseau supposée.
\par
Bornes de grille/M0 conditionnelles. Les erreurs flottantes totales
restent non majorées.
\end{frame}

\begin{frame}{Illustration : graine 20262001}
\chart{learned_curves}{0.62}
\small
Étape 3000, trois modèles sauvegardés, 257 sites.
Pour cette graine, M2 réduit les RMSE mais son score décisionnel
est légèrement plus élevé. Cette illustration n'est pas une moyenne.
\end{frame}

\begin{frame}{Les vingt paires de scores décisionnels}
\chart{paired_scores}{0.72}
\small
$D=J_{\rm phys}(\widehat p)-J_{\rm phys}(p_{\rm ref})$.
$S$ est la moyenne des trois $D$ intérieurs, sans clipping.
Vingt initialisations, toutes conservées.
\end{frame}

\begin{frame}{Gain moyen de M2 dans les deux contrastes}
\chart{paired_effects}{0.53}
\small
\begin{tabular}{lrr}
\toprule
Différence moyenne & Estimation & IC 95\,\%\\
\midrule
\Mtwo$-$\Mone & $-0.01006$ & $[-0.01460;-0.00557]$\\
\Mtwo$-$\Mthree & $-0.00940$ & $[-0.01390;-0.00503]$\\
\bottomrule
\end{tabular}
\par\vspace{0.12cm}
13 paires favorables sur 20 dans chaque comparaison.
Gain moyen, sans gain garanti pour toute graine.
\end{frame}

\begin{frame}{Erreur et mauvaise décision : panel principal}
\chart{error_decision}{0.70}
\small
M1\_16 à 3000, vingt points. Deux associations positives.
L'IC de $\Delta\rho$ inclut zéro : association supérieure
des sensibilités inconclusive.
\end{frame}

\begin{frame}{Bootstrap par blocs de graines}
\centering
\begin{tabular}{lrr}
\toprule
Statistique & Estimation & IC 95\,\%\\
\midrule
$\rho_{\rm valeur}$ &0.9413&$[0.8000;0.9882]$\\
$\rho_{\rm sens}$ &0.9481&$[0.7925;0.9958]$\\
$\Delta\rho$ &0.0068&$[-0.0568;0.0791]$\\
\bottomrule
\end{tabular}
\par\vspace{0.35cm}
\raggedright
$n=20$, $B=2000$ : une matrice commune, graines tirées avec remise,
tous les résultats du bloc conservés.
\par
Rangs moyens pour égalités exactes. Variable constante : non défini.
IC percentile interpolé, au moins 90\,\% de tirages définis.
\par
Autres conditions, étape 300 et frontières : secondaires.
Pas de pooling ou de garantie simultanée des IC.
\end{frame}

\begin{frame}{Conclusions et limites}
\begin{itemize}
\item M2 améliore le score décisionnel moyen dans les deux comparaisons prévues, dans ce benchmark.
\item L'hypothèse d'une association supérieure des sensibilités reste inconclusive.
\item Un seul cas scalaire sans bruit, vingt initialisations, référence flottante. Protocole après inspection du pilote.
\end{itemize}
\vspace{0.2cm}
Aucune nouveauté d'algorithme, causalité ou supériorité générale établie.
\par
Une extension future : transfert à un second cas physique,
avec coût des valeurs et dérivées explicitement mesuré.
Elle reste non exécutée.
\end{frame}

\begin{frame}{Sources et archives}
\small
\begin{itemize}
\item Czarnecki et al., \emph{Sobolev Training for Neural Networks},
NIPS 2017.
\item Yao et al., \emph{Derivative-Informed Fourier Neural Operator:
Universal Approximation and Applications to PDE-Constrained Optimization},
arXiv:2512.14086, 2025, révision 2026. Prépublication.
\item Alexandrov et al., \emph{A trust-region framework for managing
the use of approximation models in optimization}, 1998.
\item Elmachtoub et Grigas, \emph{Smart Predict, then Optimize},
Management Science, 2022.
\end{itemize}
Lecture bibliographique partielle ; références complètes dans le rapport.
\par
Archives : résultats 3B/3C, paramètres, contextes et exports CSV.
Aucun nouvel entraînement pour cette présentation.
\end{frame}
\end{document}
